\section{NORMAS}
El análisis y diseño estructural se rigen estrictamente por la normativa chilena vigente. Para aquellos aspectos no cubiertos por la legislación nacional, se utilizarán códigos internacionales de referencia. Las normativas aplicadas se clasifican a continuación:
  
\subsection{Cargas y Acciones sobre la Estructura}
\begin{itemize}
    \item \textbf{NCh3171:2010}: Diseño estructural -- Disposiciones generales y combinaciones de cargas \cite{NCh3171:2010}.
    \item \textbf{NCh431:2010}: Diseño estructural -- Cargas de nieve \cite{NCh431:2010}.
    \item \textbf{NCh1537:2009}: Diseño estructural -- Cargas permanentes y cargas vivas (se recomienda incluirla si mencionas cargas muertas/vivas) \cite{NCh1537:2009}.
\end{itemize}
   
\subsection{Diseño Sísmico}
\begin{itemize}
    \item \textbf{NCh433:1996 mod. 2009}: Diseño sísmico de edificios \cite{NCh433:1996}.
    \item \textbf{Decreto Supremo Nº61 (2011)}: Reglamento que fija el diseño sísmico de edificios, complementando y modificando la NCh433 \cite{DS61:2011}.
\end{itemize}

\subsection{Diseño de Elementos y Materiales}
\begin{itemize}
    \item \textbf{Decreto Supremo Nº60 (2011)}: Reglamento para el diseño y cálculo de hormigón armado \cite{DS60:2011}.
    \item \textbf{ACI 318-19}: Building Code Requirements for Structural Concrete, utilizado como referencia técnica complementaria \cite{ACI318:2019}.
    \item \textbf{NCh170:2016}: Hormigón -- Requisitos generales \cite{NCh170:2016}.
    \item \textbf{NCh204:2006}: Acero de refuerzo de hormigón -- Barras con resaltes y lisas \cite{NCh204:2006}.
\end{itemize}

\subsection{Geotecnia y Fundaciones}
\begin{itemize}
    \item \textbf{NCh1508:2014}: Geotecnia -- Estudio de mecánica de suelos \cite{NCh1508:2014}.
\end{itemize}